\section{A smaller completed direct proof-rule quartic}
\label{sec:direct-normalized-power}

The incoming direct C222 construction has a particularly useful feature:
its coefficients are ordinary printed integers, and the entire final
polynomial is explicitly materialized. It uses a finite proof-flow
certificate for twelve numerical Horn rules, followed by a CRT bounded
family and finite arithmetic representations. No Turing-machine simulation
or Jones universal polynomial occurs in this route. The incoming
fixed-base affine construction has $16036$ natural unknowns, or $27916$
unrestricted integer unknowns after its nonnegativity conversion.

We reduce these counts to $14956$ and $26188$, respectively, without
altering the twelve rules, the proof-flow encoding, or the CRT architecture.
The change concerns a redundant truncated-subtraction certificate inside
each of the $108$ calls to the power representation.

\subsection{The rank already present in the power certificate}

All variables in this subsection range over $\mathbb N=\{0,1,2,\ldots\}$.
Write the Pell residual as
\[
 \mathcal N(A;X,Y)=X^2+Y^2-A^2Y^2-1.
\]
The active finite branch of the inherited power representation contains
three facts:
\begin{equation}\label{eq:norm-rank-hypotheses}
 a>1,\qquad n\leq w,\qquad
 \mathcal N(w+1;a,wz)=0.
\end{equation}
The first follows from the outer positive-parameter guard in the Pell
representation, including when that representation uses its zero-index
branch. The second and third are explicit residuals in the power branch.

\paragraph{Rank lemma.}
Conditions~\eqref{eq:norm-rank-hypotheses} imply
\[
 w\geq1,\qquad z\geq1,\qquad a>w\geq n.
\]
Indeed, the norm equation is precisely
\[
 a^2=1+w^3(w+2)z^2.
\]
If $w=0$ or $z=0$, it gives $a^2=1$, contradicting $a>1$.
Consequently $w,z\geq1$, and
\[
 a^2=1+w^2\bigl(w(w+2)z^2\bigr)>w^2.
\]
The claimed inequality follows by monotonicity of squaring on naturals.
This is an unbounded arithmetic proof, not an inference from a finite search.

\subsection{Eliminating two semantic witnesses}

The original finite branch introduces a natural $d$ and a natural slack $s$
through the nonnegative disjunction polynomial
\begin{equation}\label{eq:old-truncated-subtraction}
 (a-d-n)^2\bigl((n-a-s)^2+d^2\bigr)=0.
\end{equation}
It expresses $d=\max(a-n,0)$. Under the rank lemma, the second factor
cannot vanish, since $a>n$ makes $n-a-s\ne0$. Thus $d=a-n$.
The same branch contains the congruence, written with natural quotient
witnesses $q_0,q_1$,
\[
 x+tq_0=yd+m+tq_1.
\]
Substitute $d=a-n$ and move $yn$ to the left. The resulting equation is
\begin{equation}\label{eq:translated-power-congruence}
 x+yn+tq_0=ya+m+tq_1.
\end{equation}
Both sides are natural arithmetic expressions. The new compiler deletes
$d$, $s$, and~\eqref{eq:old-truncated-subtraction}, and replaces the old
congruence by~\eqref{eq:translated-power-congruence}. Every other byte of
the original compiler is retained; the delivered delta checker verifies this.

\paragraph{Soundness and completeness, including inactive branches.}
When the finite branch is selected, its residuals form a sum of
nonnegative expressions, so each vanishes. The rank lemma applies to the
new branch as well as the old one. An old witness therefore restricts to a
new witness. Conversely, a new finite-branch witness extends to an old one
by setting $d=a-n$ and $s=0$. The first factor in
\eqref{eq:old-truncated-subtraction} is then zero, and the two congruences
are the same integer-polynomial equation.

The inherited power representation separately covers exponent zero and
zero base with positive exponent. Neither branch is changed. When one of
these branches is selected, the finite expression is underneath a zero
factor in the disjunction, so it imposes no condition. The erased witnesses
can then be restored as $d=s=0$. All deleted or changed natural arithmetic
wires can be recomputed from their inputs. There is no requirement that an
unselected subtraction or Pell branch vanish. This observation is essential
to the two-way existential equivalence of the complete power graph.

\subsection{Composition into one literal polynomial}

Replacing an existential arithmetic graph by an equivalent graph preserves
every enclosing binomial, factorial and bounded-CRT representation. The
unchanged global construction uses $108$ power graphs, $33$ binomial graphs,
$3$ factorial graphs and $72$ positive-divisor Euclidean-division graphs.
All twelve local Horn branches and both global instantiations of the local
row remain unchanged.

Let $\rho_1,\ldots,\rho_{11502}$ be the explicit quadratic residuals in
\texttt{residuals\_affine.json}, in $14956$ natural variables. The natural
equation is
\[
 P_{\mathbb N}(X)=\sum_{j=1}^{11502}\rho_j(X)^2=0.
\]
The residuals are literal finite integer polynomials, supplied in sparse
form, and their complete expansion is \texttt{equation\_affine.json}.
There are $3744$ semantic inputs and $11212$ retained nonlinear wires.

To obtain the unrestricted-integer equation, introduce three fresh integer
variables $u_i,v_i,w_i$ for each semantic input $X_i$, $0\leq i<3744$, and set
\[
 P_{\mathbb Z}
 =P_{\mathbb N}
  +\sum_{i=0}^{3743}
    \bigl(X_i-u_i^2-v_i^2-w_i(w_i+1)\bigr)^2.
\]
The polynomial has $14956+3\cdot3744=26188$ integer variables and total
degree exactly four. Every natural is of the form
$u^2+v^2+w(w+1)$ with integer $u,v,w$: apply the three-square theorem to
$4X+1$, whose residue modulo eight is one or five, and separate the single
odd square from the two even squares. Conversely, the displayed expression
is nonnegative for every integer triple. All other wires are forced to be
nonnegative by the unchanged topological equations and nonnegative affine
forms, so they need no additional guards.

The complete integer expansion is \texttt{equation\_integer.json}; it has
$724768$ nonzero monomials and maximum absolute coefficient
$320307372009864$. It contains no variable-exponent, factorial, binomial,
bounded-quantifier or symbolic-coefficient operation.

\subsection{Logical calibration and evidence boundary}

The unchanged proof-flow certificate is equivalent to a finite derivation of
the distinguished target in C222. The pinned logical interfaces are
\texttt{KEndpoint.c222\_iff\_zf\_refutation} and
\texttt{ConstructibleChoice.zf\_refutation\_iff\_zfc\_refutation}.
The fixed Horn coefficients are independently matched to
\texttt{Code254.lean} and \texttt{KCode222.lean}; the names of those files
do not introduce a Turing-machine simulation into this construction.
The inherited uniform CRT/Pell proof, the normalization proved above, and
the finite circuit lifting yield ordinary mathematical proofs formalizable
in ZFC of
\[
 \operatorname{ZFC}\vdash
 \bigl(\neg\exists X\in\mathbb N^{14956}:P_{\mathbb N}(X)=0\bigr)
 \;\longleftrightarrow\;\operatorname{Con}(\operatorname{ZFC}),
\]
and the corresponding statement for $P_{\mathbb Z}$ over
$\mathbb Z^{26188}$. No consistency assumption is used in either direction.

The exact finite audits check all arithmetic-interface polynomials,
compositions, Horn branches, source substitutions, gate residuals,
nonnegative affine eliminations and full expanded polynomial identities.
The original arithmetic and logical source files are supplied with immutable
hashes. These checks do not amount to a newly compiled Lean or Rocq proof of
the complete construction. The claimed result is an ordinary proof with
fully materialized and exactly audited polynomial artifacts.

This is a strict reduction of the incoming completed direct construction.
It is not a global fewest-variable result, and it does not dominate the
repository's native low-variable equations. The elementary rank argument is
not claimed to be historically new; its useful application here is the
removal of a redundant branch certificate from every power call.
